\begin{abstract}
  This paper introduces Norm-Equalized Task Arithmetic (NETA), a training-free and parameter-free closed-form method for multi-task model merging. While standard task arithmetic directly sums fine-tuned weight updates, it is highly sensitive to disparity in the norms of individual task vectors, allowing high-magnitude tasks to dominate the representation space and destructively interfere with others. NETA addresses this imbalance analytically by equalizing the Frobenius norms of task vectors at each layer before merging, ensuring isotropic contribution across all task experts. In visual classification experiments with CLIP ViT-B/32, NETA outperforms standard task arithmetic zero-shot across two of four target tasks. While NETA's role as an isotropic regularizer slightly curtails peak performance on dominant tasks (resulting in a marginal drop in multi-task average accuracy from $87.76\%$ to $87.17\%$), it successfully enforces representational fairness and prevents task suppression. We also present a continuous $\alpha$-relaxation framework to smoothly balance representation fairness and peak task performance. Crucially, we demonstrate that state-of-the-art test-time weight optimization techniques suffer from an ``Overfitting-Optimizer Paradox''---where joint entropy minimization overfits to easy, low-entropy tasks and suppresses complex ones---whereas NETA's simple, analytical design achieves stable performance and high robustness without any calibration data, learning rates, or backpropagation.
\end{abstract}
